\% !TEX root = ../../buch.tex \% !TEX encoding = UTF-8 \%
\textbackslash section\{Die Padé-Approximation\}
\textbackslash label\{pade:section:pade\}
\textbackslash kopfrechts\{Padé-Approximation\}

Im vorherigen Abschnitt wurde gezeigt, dass Taylor-Polynome eine
Funktion lediglich in der Umgebung ihres Entwicklungspunktes gut
approximieren. Obwohl die Taylor-Reihe wichtige Informationen über das
lokale Verhalten einer Funktion liefert, nimmt ihre Genauigkeit mit
zunehmendem Abstand vom Entwicklungspunkt häufig deutlich ab. Daraus
ergibt sich unmittelbar die Frage, ob sich dieselben Informationen der
Taylor-Reihe auf eine andere Weise nutzen lassen.

Genau diese Fragestellung führte Henri Padé Ende des
19.\textasciitilde Jahrhunderts zur Entwicklung seines
Approximationsverfahrens. Anstatt die Koeffizienten der Taylor-Reihe
direkt zu einem Polynom zusammenzufassen, werden sie verwendet, um eine
rationale Funktion zu bestimmen. Unter einer rationalen Funktion
versteht man den Quotienten zweier Polynome. Durch den zusätzlichen
Nenner besitzt eine solche Funktion mehr Freiheitsgrade als ein
einzelnes Polynom und kann den Verlauf vieler Funktionen genauer
nachbilden.

Allgemein besitzt eine Padé-Approximation vom Typ
\textbackslash({[}m/n{]}\textbackslash) die Form
\textbackslash begin\{equation\} R\_\{m,n\}(x) =
\textbackslash frac\{P\_m(x)\}\{Q\_n(x)\} =
\textbackslash frac\{a\_0+a\_1x+\textbackslash cdots+a\_mx\^{}m\}
\{1+b\_1x+\textbackslash cdots+b\_nx\^{}n\}.
\textbackslash label\{pade:eq:allgemein\} \textbackslash end\{equation\}

Dabei bezeichnet \textbackslash(m\textbackslash) den Grad des
Zählerpolynoms und \textbackslash(n\textbackslash) den Grad des
Nennerpolynoms. Die Koeffizienten \textbackslash(a\_i\textbackslash) und
\textbackslash(b\_i\textbackslash) sind zunächst unbekannt und müssen
mithilfe der Taylor-Reihe bestimmt werden. Der konstante Term des
Nennerpolynoms wird auf den Wert \textbackslash(1\textbackslash)
normiert. Diese Normierung verändert die Approximation nicht, sorgt
jedoch dafür, dass jede Padé-Approximation eindeutig dargestellt wird.
Würde auch dieser Koeffizient als Unbekannte betrachtet, könnten
dieselben rationalen Funktionen durch beliebig viele verschiedene
Koeffizientensätze beschrieben werden. Zudem müssen die Koeffizienten
\textbackslash(a\_m\textbackslash) und
\textbackslash(b\_n\textbackslash) von null verschieden sein.
Andernfalls wäre der tatsächliche Grad des Zähler- beziehungsweise
Nennerpolynoms kleiner und es würde sich nicht mehr um eine echte
\textbackslash({[}m/n{]}\textbackslash)-Approximation handeln.

Das Ziel besteht darin, die unbekannten Koeffizienten so zu bestimmen,
dass die Reihenentwicklung der Padé-Approximation möglichst viele
Anfangsterme der Taylor-Reihe reproduziert. Die Taylor-Reihe liefert
somit sämtliche Informationen, welche zur Konstruktion der
Padé-Approximation benötigt werden. Zusätzliche Ableitungen oder weitere
Funktionswerte sind nicht erforderlich.

Allgemein gilt für eine Padé-Approximation vom Typ
\textbackslash({[}m/n{]}\textbackslash), dass sich ihre
Reihenentwicklung möglichst lange mit der Taylor-Reihe der
ursprünglichen Funktion decken soll. Mathematisch lässt sich diese
Forderung durch \textbackslash begin\{equation\} f(x)-R\_\{m,n\}(x) =
O\textbackslash!\textbackslash left(x\^{}\{m+n+1\}\textbackslash right)
\textbackslash label\{pade:eq:ordnung\} \textbackslash end\{equation\}
ausdrücken. Die Landau-Notation \textbackslash(O\textbackslash)
beschreibt hierbei die Ordnung des verbleibenden Approximationsfehlers.
Die Gleichung bedeutet, dass sich die Taylor-Reihe der Funktion und die
Reihenentwicklung der Padé-Approximation erst ab der Potenz
\textbackslash(x\^{}\{m+n+1\}\textbackslash) unterscheiden. Alle
vorhergehenden Koeffizienten stimmen exakt überein.

Zur Veranschaulichung wird im Folgenden die Exponentialfunktion
betrachtet. Für diese Funktion wird Schritt für Schritt eine
Padé-Approximation vom Typ \textbackslash({[}1/1{]}\textbackslash)
hergeleitet. Dieses Beispiel zeigt das grundlegende Prinzip des
Verfahrens und bildet gleichzeitig die Grundlage für Approximationen
höherer Ordnung.

\textbackslash subsection\{Herleitung der Padé-{[}1/1{]}-Approximation\}
\textbackslash label\{pade:subsection:pade11\}

Zur Veranschaulichung wird die Exponentialfunktion
\textbackslash(f(x)=e\^{}x\textbackslash) betrachtet. Gesucht wird eine
Padé-Approximation vom Typ \textbackslash({[}1/1{]}\textbackslash). Dies
bedeutet, dass sowohl der Zähler als auch der Nenner aus einem Polynom
ersten Grades bestehen.

Der allgemeine Ansatz lautet \textbackslash begin\{equation\}
R\_\{1,1\}(x) = \textbackslash frac\{a\_0+a\_1x\}\{1+b\_1x\}.
\textbackslash label\{pade:eq:pade11\_ansatz\}
\textbackslash end\{equation\}

Da der Zähler zwei unbekannte Koeffizienten und der Nenner einen
weiteren unbekannten Koeffizienten enthält, müssen insgesamt drei
Unbekannte bestimmt werden. Dementsprechend werden auch drei unabhängige
Gleichungen benötigt.

Diese Gleichungen erhält man durch den Vergleich mit der Taylor-Reihe
der Exponentialfunktion. Aus
Abschnitt\textasciitilde\textbackslash ref\{pade:section:taylor\} ist
bekannt, dass diese lautet \textbackslash begin\{equation\} e\^{}x = 1 +
x + \textbackslash frac\{x\^{}2\}\{2\} +
\textbackslash frac\{x\^{}3\}\{6\} + \textbackslash frac\{x\^{}4\}\{24\}
+ \textbackslash cdots. \textbackslash label\{pade:eq:exp\_taylor\}
\textbackslash end\{equation\}

Ein direkter Vergleich der Koeffizienten ist zunächst nicht möglich, da
die Padé-Approximation als Bruch vorliegt. Deshalb muss zunächst der
Nenner entwickelt werden.

Hierzu wird die geometrische Reihe verwendet. Für
\textbackslash(\textbar y\textbar\textless1\textbackslash) gilt
\textbackslash begin\{equation\} \textbackslash frac\{1\}\{1+y\} =
1-y+y\^{}2-y\^{}3+y\^{}4-\textbackslash cdots.
\textbackslash label\{pade:eq:geometrisch\}
\textbackslash end\{equation\}

Setzt man \textbackslash(y=b\_1x\textbackslash), so erhält man
\textbackslash begin\{equation\} \textbackslash frac\{1\}\{1+b\_1x\} = 1
- b\_1x + b\_1\^{}2x\^{}2 - b\_1\^{}3x\^{}3 +\textbackslash cdots.
\textbackslash label\{pade:eq:nennerentwicklung\}
\textbackslash end\{equation\}

Damit kann auch die Padé-Approximation als Potenzreihe dargestellt
werden: \textbackslash begin\{equation\} R\_\{1,1\}(x) = (a\_0+a\_1x)
\textbackslash left( 1 - b\_1x + b\_1\^{}2x\^{}2 - b\_1\^{}3x\^{}3
+\textbackslash cdots \textbackslash right).
\textbackslash label\{pade:eq:potenzreihe\}
\textbackslash end\{equation\}

Im nächsten Schritt wird dieses Produkt ausmultipliziert. Dabei wird
zunächst jeder Summand des Zählerpolynoms mit jedem Summanden der
Potenzreihe multipliziert. Anschliessend werden alle Terme gleicher
Ordnung zusammengefasst. Dadurch erhält man
\textbackslash begin\{equation\} R\_\{1,1\}(x) = a\_0 + (a\_1-a\_0b\_1)x
+ (a\_0b\_1\^{}2-a\_1b\_1)x\^{}2 + (a\_1b\_1\^{}2-a\_0b\_1\^{}3)x\^{}3
+\textbackslash cdots. \textbackslash label\{pade:eq:entwicklung\}
\textbackslash end\{equation\}

Nun liegen sowohl die Exponentialfunktion als auch die
Padé-Approximation in Form einer Potenzreihe vor. Dadurch können die
Koeffizienten gleicher Potenzen direkt miteinander verglichen werden.
Aus dem Koeffizientenvergleich ergeben sich die drei Gleichungen
\textbackslash begin\{equation\} \textbackslash begin\{aligned\} a\_0
\&= 1,\textbackslash\textbackslash{} a\_1-a\_0b\_1 \&=
1,\textbackslash\textbackslash{} a\_0b\_1\^{}2-a\_1b\_1 \&=
\textbackslash frac\{1\}\{2\}. \textbackslash end\{aligned\}
\textbackslash label\{pade:eq:gleichungssystem\}
\textbackslash end\{equation\}

Aus der ersten Gleichung folgt unmittelbar
\textbackslash(a\_0=1\textbackslash). Damit reduziert sich das
Gleichungssystem auf \textbackslash begin\{equation\}
\textbackslash begin\{aligned\} a\_1-b\_1 \&=
1,\textbackslash\textbackslash{} b\_1\^{}2-a\_1b\_1 \&=
\textbackslash frac\{1\}\{2\}. \textbackslash end\{aligned\}
\textbackslash label\{pade:eq:gleichungssystem\_reduziert\}
\textbackslash end\{equation\}

Die linke Seite der zweiten Gleichung lässt sich als \textbackslash{[}
b\_1(b\_1-a\_1) \textbackslash{]} schreiben. Aus der ersten Gleichung
folgt \textbackslash(a\_1-b\_1=1\textbackslash) beziehungsweise
\textbackslash(b\_1-a\_1=-1\textbackslash). Somit ergibt sich
\textbackslash{[} -b\_1=\textbackslash frac\{1\}\{2\}, \textbackslash{]}
woraus unmittelbar \textbackslash begin\{equation\}
b\_1=-\textbackslash frac\{1\}\{2\} \textbackslash label\{pade:eq:b1\}
\textbackslash end\{equation\} folgt.

Durch Einsetzen von
\textbackslash(b\_1=-\textbackslash frac\{1\}\{2\}\textbackslash) in die
erste Gleichung des reduzierten Systems erhält man
\textbackslash begin\{equation\} a\_1 = 1-\textbackslash frac\{1\}\{2\}
= \textbackslash frac\{1\}\{2\}. \textbackslash label\{pade:eq:a1final\}
\textbackslash end\{equation\}

Damit sind sämtliche unbekannten Koeffizienten bestimmt. Durch Einsetzen
in
Gleichung\textasciitilde\textbackslash eqref\{pade:eq:pade11\_ansatz\}
ergibt sich die Padé-Approximation \textbackslash begin\{equation\}
R\_\{1,1\}(x) = \textbackslash frac\{1+\textbackslash frac\{x\}\{2\}\}
\{1-\textbackslash frac\{x\}\{2\}\}.
\textbackslash label\{pade:eq:pade11\} \textbackslash end\{equation\}

Dieses Ergebnis stellt die Padé-{[}1/1{]}-Approximation der
Exponentialfunktion dar. Obwohl zu ihrer Herleitung ausschliesslich die
ersten Koeffizienten der Taylor-Reihe verwendet wurden, entsteht keine
Polynomfunktion, sondern eine rationale Funktion. Gerade diese
unterschiedliche mathematische Struktur führt häufig dazu, dass die
Padé-Approximation den Verlauf einer Funktion genauer beschreibt als ein
Taylor-Polynom gleicher Ordnung.

Um diesen Unterschied zu verdeutlichen, wird im folgenden Abschnitt die
Padé-{[}1/1{]}-Approximation anhand eines konkreten Zahlenbeispiels mit
dem Taylor-Polynom zweiter Ordnung verglichen.
